% Soluzioni degli esercizi del Modulo 7 (generato da separa_soluzioni.py)
\chapter{Soluzioni -- Modulo 7 -- Algoritmi di Ordinamento e Ricerca}
\textbf{Esercizio 7.1}
\codice{codice/m07/es01.cpp}
\uscita{codice/m07/output/es01.txt}

\textbf{Esercizio 7.2}
\codice{codice/m07/es02.cpp}
\uscita{codice/m07/output/es02.txt}

\textbf{Esercizio 7.3} --- Nel codice ASCII le maiuscole (65--90) vengono prima delle
minuscole (97--122), quindi \q{bianchi} va dopo \q{Verdi}.
\codice{codice/m07/es03.cpp}
\uscita{codice/m07/output/es03.txt}

\textbf{Esercizio 7.4} --- Il risultato si salva in una variabile \emph{prima} di
stamparlo: se la chiamata fosse dentro lo stesso \cc|cout|, la scritta
\texttt{-> indice} comparirebbe prima dei passi.
\codice{codice/m07/es04.cpp}
\uscita{codice/m07/output/es04.txt}

\textbf{Esercizio 7.5}
\codice{codice/m07/es05.cpp}
\uscita{codice/m07/output/es05.txt}

\textbf{Esercizio 7.6} --- Con \cc|<|, a parità di voto verrebbe preso prima
l'elemento della metà destra, e l'ordine alfabetico si romperebbe.
\codice{codice/m07/es06.cpp}
\uscita{codice/m07/output/es06.txt}

\textbf{Esercizio 7.7}
\codice{codice/m07/es07.cpp}
\uscita{codice/m07/output/es07.txt}

\textbf{Esercizio 7.8} --- $O(n + k)$, con $k$ = numero di valori possibili (13).
\codice{codice/m07/es08.cpp}
\uscita{codice/m07/output/es08.txt}

\textbf{Esercizio 7.9}
\codice{codice/m07/es09.cpp}
\uscita{codice/m07/output/es09.txt}

\textbf{Esercizio 7.10}
\codice{codice/m07/es10.cpp}
\uscita{codice/m07/output/es10.txt}
